\documentclass[12pt,a4paper]{article}
\usepackage{amsmath}
\usepackage{amssymb}
\usepackage{graphicx}
\usepackage{subcaption}
\usepackage{enumitem}
\usepackage{listings}
\usepackage{hyperref}
\usepackage{ulem}
\usepackage[table]{xcolor}
\usepackage{amsthm}
\usepackage{tcolorbox}
\usepackage{array}
\usepackage{mdframed}
\usepackage{cancel}
\usepackage[margin=2.5cm]{geometry}

\setlength{\parindent}{0pt}
\setlength{\parskip}{0pt}
\setcounter{MaxMatrixCols}{11}

\begin{document}
\begin{titlepage}
\centering
\vspace*{\fill}
{\Huge Algebra}\\[0.5cm]
{\Huge Home Work 8}\\[1.5cm]
{\Large Student: 26-712-224}\\[0.5cm]
{\Large \today}\par
\vspace{2cm}
{\normalsize ``I wrote these solutions fully myself und would like to receive feedback on these''\par}
\vspace*{\fill}
\end{titlepage}
% ============================================================
% Question A
% ============================================================
\noindent For a finite group $G$ and a prime number $p$, we denote by $s_p(G)$ the cardinality of the set $\mathrm{Syl}_p(G)$ of $p$-Sylow subgroups of $G$.\\
\\
\noindent\textbf{A)} Let $G$ be a finite group of order $200$. Show that $G$ is not simple. (\textit{Hint:} Determine $s_5(G)$.)
\newpage
% ============================================================
% Question B
% ============================================================
\noindent\textbf{B)} Let $p$ and $q$ be prime numbers and $G$ a finite group of order $pq$. Show that $G$ is not simple. (\textit{Hint:} When $p<q$, determine $s_q(G)$.)
\newpage
% ============================================================
% Question C
% ============================================================
\noindent\textbf{C)} We want to show that there is no simple group of order $400$. For this, we assume that there is a finite simple group $G$ of order $400$ and we obtain a contradiction as follows.
\begin{enumerate}
\item[(1)] Show that $s_5(G)=16$ and $s_2(G)\in\{5,25\}$.
\item[(2)] Let $S\subset G$ be a $5$-Sylow subgroup. Show that $S$ is abelian. (\textit{Hint:} Use Exercise \textbf{A)} of Sheet 6.)
\item[(3)] Let $S_1$ and $S_2$ be two distinct $5$-Sylow subgroups of $G$. We want to show that $S_1\cap S_2=\{1\}$. We argue by contradiction and assume that $Z=S_1\cap S_2$ is different form $\{1\}$.
\begin{enumerate}
\item[(a)] Show that $|Z|=5$.
\item[(b)] Let $H=N_G(Z)$ be the normalizer of $Z$. Show that $S_1$ and $S_2$ are contained in $H$. (\textit{Hint:} Use (2).)
\item[(c)] Show that $H\neq G$. (Use that $G$ is simple.) Deduce that $|H|\in\{50,100,200\}$.
\item[(d)] Use (c) to show that $s_5(H)=1$. Deduce a contradiction.
\end{enumerate}
\item[(4)] Let $P_5\subset G$ be the set of elements of order in $\{5,25\}$. Use (1) and (3) to show that $\mathrm{card}(P_5)=384$.
\item[(5)] Let $P_2\subset G$ be the set of elements of order in $\{2,4,8,16\}$. Use (1) to show that $\mathrm{card}(P_2)\geq 23$.
\item[(6)] Use (4) and (5) to obtain a contradiction.
\end{enumerate}
\newpage
% ============================================================
% Question D
% ============================================================
\noindent\textbf{D)} Consider the permutation $\sigma\in\Sigma_{11}$ given by:
\begin{align*}
\sigma=
\begin{pmatrix}
1&2&3&4&5&6&7&8&9&10&11\\
3&5&7&6&11&4&1&2&8&10&9
\end{pmatrix}.
\end{align*}
\begin{enumerate}
\item[(1)] Find the decomposition into cycles of $\sigma$.
\item[(2)] Compute the order of $\sigma$.
\item[(3)] Compute the signature of $\sigma$.
\end{enumerate}

\vspace{1cm}
\noindent
Answer:\\
\\
(1) We can represent $\sigma$ as the following cycles:
\begin{align*}
    \sigma &= (1 \; 3 \; 7)(2 \; 5 \; 11 \; 9 \; 8)(4 \; 6) \\
        &= c_1 \circ c_2 \circ c_3
\end{align*}
(2) The order of $\sigma$ is:
\begin{align*}
 \text{ord}(\sigma) &= \text{lcm}(\text{ord}(c_1),\; \text{ord}(c_2),\; \text{ord}(c_3)) \\
   &= \text{lcm}(3, 5, 2) \\
   &= 30
\end{align*}
(3) The signature of $\sigma$ is:
\begin{align*}
 \text{sgn}(\sigma) &= \text{sgn}(c_1) . \text{sgn}(c_2) . \text{sgn}(c_3)\\
  &= (-1)^{\text{ord}(c_1)-1} . (-1)^{\text{ord}(c_2)-1} . (-1)^{\text{ord}(c_3)-1} \\
  &= (-1)^2 . (-1)^4 . (-1)^1 \\
  &= -1
\end{align*}

\end{document}